% ============================================================================
%  MATEMÁTICAS OPERATIVAS — Sesión 8
%  Semana 8: 28 de septiembre al 3 de octubre de 2026
%  2.2 Radicación y sus propiedades
%  Duración prevista: 3 horas
% ============================================================================
\documentclass[aspectratio=169,10pt]{beamer}

\usepackage[utf8]{inputenc}
\IfFileExists{spanish.ldf}{\usepackage[spanish,es-noshorthands]{babel}}{}
\usepackage{amsmath,amssymb}
\usepackage{booktabs}
\usepackage{tikz}

\usetheme{Uniminuto}

\uniasignatura{Matemáticas Operativas}
\uniprofesor{Andrés Felipe Puerta Velez}
\unicorreo{andres.puerta-v@uniminuto.edu.co}
\uninrc{CBASBL011}
\unisemana{Sesión 8 — Semana 8 · 28 de septiembre al 3 de octubre de 2026}

\newcommand{\listatiny}{\setbeamerfont{itemize/enumerate body}{size=\tiny}\tiny}

\begin{document}

\portada

\begin{frame}{Orden del día}
  \begin{enumerate}
    \item Revisión de los ejercicios propuestos de la sesión 7
    \item \azul{Tema 2.2: la radicación como potencia fraccionaria}
    \item Propiedades de los radicales
    \item Simplificar radicales y extraer factores
    \item Sumar, restar, multiplicar y dividir radicales
    \item Taller principal, errores frecuentes y cierre
  \end{enumerate}

  \vspace{0.25cm}
  {\footnotesize\color{uniGris} Sesión de 3 horas. Todo el tema se apoya en una
  sola idea: \destacar{un radical es una potencia con exponente fraccionario}.
  Quien la asuma no tiene que memorizar propiedades nuevas.}
\end{frame}

% ===========================================================================
\seccion{2.2 Radicación y sus propiedades}
% ===========================================================================

\begin{frame}{La idea que resuelve el tema}
  \scriptsize
  \begin{block}{Radical y potencia son lo mismo}
    \scriptsize
    \[ \sqrt[n]{a^{m}} = a^{\frac{m}{n}} \]
    El \azul{índice} de la raíz va al \azul{denominador} del exponente; el
    exponente del radicando, al numerador.
  \end{block}

  \vspace{0.1cm}
  \begin{center}
  \begin{tabular}{@{}llll@{}}
    \toprule
    \textbf{Radical} & \textbf{Potencia} & \textbf{Valor} &
    \textbf{Cómo se piensa} \\
    \midrule
    $\sqrt[3]{8^{2}}$ & $8^{2/3}$ & $4$ &
      La raíz cúbica de $8$ es $2$; $2^{2}=4$ \\
    $\sqrt[4]{16^{3}}$ & $16^{3/4}$ & $8$ &
      $\sqrt[4]{16}=2$; $2^{3}=8$ \\
    $\sqrt[5]{32^{3}}$ & $32^{3/5}$ & $8$ &
      $\sqrt[5]{32}=2$; $2^{3}=8$ \\
    $\dfrac{1}{\sqrt[3]{27^{2}}}$ & $27^{-2/3}$ & $\tfrac{1}{9}$ &
      $\sqrt[3]{27}=3$; $3^{-2}=\tfrac{1}{9}$ \\
    \bottomrule
  \end{tabular}
  \end{center}

  \vspace{0.08cm}
  {\tiny\color{uniGris} \destacar{Truco de cálculo:} saque \azul{primero la
  raíz} y después eleve. $16^{3/4}$ es más fácil como $(\sqrt[4]{16})^{3}=8$
  que como $\sqrt[4]{4096}$.}
\end{frame}

\begin{frame}{Propiedades de los radicales}
  \scriptsize
  \begin{center}
  \begin{tabular}{@{}p{4.2cm}p{3.8cm}p{4.0cm}@{}}
    \toprule
    \textbf{Propiedad} & \textbf{Regla} & \textbf{Ejemplo} \\
    \midrule
    Raíz de un producto &
      $\sqrt[n]{ab} = \sqrt[n]{a}\cdot\sqrt[n]{b}$ &
      $\sqrt{5}\cdot\sqrt{20} = \sqrt{100} = 10$ \\
    \addlinespace
    Raíz de un cociente &
      $\sqrt[n]{\dfrac{a}{b}}
       = \dfrac{\sqrt[n]{a}}{\sqrt[n]{b}}$ &
      $\dfrac{\sqrt{98}}{\sqrt{2}} = \sqrt{49} = 7$ \\
    \addlinespace
    Potencia de un radical &
      $\left(\sqrt[n]{a}\right)^{m} = \sqrt[n]{a^{m}}$ &
      $(\sqrt{3})^{4} = \sqrt{81} = 9$ \\
    \addlinespace
    Raíz de una raíz &
      $\sqrt[m]{\sqrt[n]{a}} = \sqrt[mn]{a}$ &
      $\sqrt{\sqrt[3]{a}} = \sqrt[6]{a}$ \\
    \bottomrule
  \end{tabular}
  \end{center}

  \vspace{0.1cm}
  \begin{alertblock}{La propiedad que NO existe}
    \tiny
    \[ \sqrt{a+b} \;\neq\; \sqrt{a}+\sqrt{b} \]
    Compruebe: $\sqrt{9+16}=\sqrt{25}=5$, pero $\sqrt{9}+\sqrt{16}=3+4=7$.
    \destacar{La raíz se reparte en productos y cocientes, nunca en sumas.}
    Es exactamente el mismo patrón que con las potencias.
  \end{alertblock}
\end{frame}

\begin{frame}{Simplificar: extraer factores del radical}
  \scriptsize
  \begin{block}{El método}
    \tiny
    \textbf{1.} Factorice el radicando en primos. \quad
    \textbf{2.} Agrupe los factores en bloques del tamaño del \azul{índice}.
    \quad
    \textbf{3.} Cada bloque completo \destacar{sale} con un solo factor; lo
    que sobra se queda dentro.
  \end{block}

  \vspace{0.1cm}
  \begin{columns}[T,onlytextwidth]
    \begin{column}{0.48\textwidth}
      \begin{exampleblock}{Índice 2}
        \tiny
        \[ \sqrt{72} = \sqrt{2^{3}\cdot 3^{2}}
           = \sqrt{2^{2}\cdot 3^{2}\cdot 2} \]
        \[ = 2\cdot 3\sqrt{2} = \destacar{6\sqrt{2}} \]
        \azul{Comprobación:} $6\sqrt{2}\approx 8{,}485$ y
        $\sqrt{72}\approx 8{,}485$ \checkmark
      \end{exampleblock}
    \end{column}
    \begin{column}{0.48\textwidth}
      \begin{exampleblock}{Índice 3}
        \tiny
        \[ \sqrt[3]{54} = \sqrt[3]{2\cdot 3^{3}}
           = \destacar{3\sqrt[3]{2}} \]
        \[ \sqrt[3]{128} = \sqrt[3]{2^{7}}
           = \destacar{4\sqrt[3]{2}} \]
        {\color{uniGris}$2^{7}=2^{6}\cdot 2$, y $2^{6}$ son dos bloques de
        tres.}
      \end{exampleblock}
    \end{column}
  \end{columns}

  \vspace{0.08cm}
  {\tiny\color{uniGris} \azul{Con letras es igual:}
  $\sqrt{50x^{3}} = \sqrt{5^{2}\cdot 2\cdot x^{2}\cdot x}
  = \destacar{5x\sqrt{2x}}$ \ (con $x\ge 0$). El exponente se divide entre el
  índice: cociente afuera, resto adentro.}
\end{frame}

\begin{frame}{Ejercicio 1 — simplificar radicales}
  \scriptsize
  \begin{alertblock}{Individual, 10 minutos}
    \tiny
    \textbf{1.} $\sqrt{200}$ \qquad
    \textbf{2.} $\sqrt{48}$ \qquad
    \textbf{3.} $\sqrt[4]{80}$ \qquad
    \textbf{4.} $\sqrt{18x^{4}y^{6}}$ \qquad
    \textbf{5.} $\sqrt[3]{24x^{6}}$
  \end{alertblock}

  \vspace{0.06cm}
  \begin{block}{Solución}
    \tiny
    \textbf{1.} $\sqrt{2^{3}5^{2}} = \destacar{10\sqrt{2}}$ \qquad
    \textbf{2.} $\sqrt{2^{4}\cdot 3} = \destacar{4\sqrt{3}}$ \qquad
    \textbf{3.} $\sqrt[4]{2^{4}\cdot 5} = \destacar{2\sqrt[4]{5}}$

    \vspace{0.08cm}
    \textbf{4.} $\sqrt{3^{2}\cdot 2\cdot x^{4}y^{6}}
    = \destacar{3x^{2}y^{3}\sqrt{2}}$ \qquad
    \textbf{5.} $\sqrt[3]{2^{3}\cdot 3\cdot x^{6}}
    = \destacar{2x^{2}\sqrt[3]{3}}$
  \end{block}

  \vspace{0.05cm}
  {\tiny\color{uniGris} \azul{Atajo para las letras:} divida el exponente
  entre el índice. En el \destacar{4}, $x^{4}$ con índice $2$ da $x^{2}$
  afuera y nada adentro; $y^{6}$ da $y^{3}$. En el \destacar{5}, $x^{6}$ con
  índice $3$ da $x^{2}$.}
\end{frame}

\begin{frame}{Sumar y restar radicales}
  \scriptsize
  \begin{block}{Solo se suman los radicales semejantes}
    \tiny
    Semejantes significa \destacar{mismo índice y mismo radicando}. Se suman
    los coeficientes y el radical se conserva, igual que con términos
    algebraicos: $3\sqrt{2}+5\sqrt{2} = 8\sqrt{2}$, como $3z+5z=8z$.
  \end{block}

  \vspace{0.1cm}
  \begin{exampleblock}{El paso que casi nadie da: simplificar primero}
    \tiny
    \[ \sqrt{12}+\sqrt{27} \]
    A primera vista no son semejantes. Pero al simplificar:
    \[ 2\sqrt{3}+3\sqrt{3} = \destacar{5\sqrt{3}} \]
    Otro caso:
    \[ 3\sqrt{8}-\sqrt{50} = 6\sqrt{2}-5\sqrt{2} = \destacar{\sqrt{2}} \]
  \end{exampleblock}

  \vspace{0.0cm}
  \begin{alertblock}{La regla de oro del tema}
    \tiny
    \destacar{Antes de decidir si dos radicales se pueden sumar,
    simplifíquelos.} Muchos radicales que parecen distintos resultan
    semejantes; y si tras simplificar siguen sin serlo, la suma
    \azul{se deja indicada}: $\sqrt{2}+\sqrt{3}$ no se puede reducir más.
  \end{alertblock}
\end{frame}

\begin{frame}{Ejercicio 2 — operar con radicales}
  \scriptsize
  \begin{alertblock}{En parejas, 12 minutos}
    \tiny
    \textbf{1.} $2\sqrt{45}+3\sqrt{20}-\sqrt{80}$ \qquad
    \textbf{2.} $\sqrt{6}\cdot\sqrt{15}$ \qquad
    \textbf{3.} $\dfrac{\sqrt{162}}{\sqrt{2}}$ \qquad
    \textbf{4.} $\sqrt{x}\cdot\sqrt[3]{x}$ \qquad
    \textbf{5.} $\left(\sqrt[3]{5}\right)^{6}$
  \end{alertblock}

  \vspace{0.05cm}
  \begin{block}{Solución}
    \tiny
    \textbf{1.} $2(3\sqrt5)+3(2\sqrt5)-4\sqrt5
    = 6\sqrt5+6\sqrt5-4\sqrt5 = \destacar{8\sqrt{5}}$

    \vspace{0.08cm}
    \textbf{2.} $\sqrt{90} = \destacar{3\sqrt{10}}$ \qquad
    \textbf{3.} $\sqrt{81} = \destacar{9}$

    \vspace{0.08cm}
    \textbf{4.} $x^{1/2}\cdot x^{1/3} = x^{1/2+1/3}
    = x^{5/6} = \destacar{\sqrt[6]{x^{5}}}$ \qquad
    \textbf{5.} $5^{6/3} = 5^{2} = \destacar{25}$
  \end{block}

  \vspace{0.05cm}
  {\tiny\color{uniGris} El punto \destacar{4} es la razón de ser de la
  notación fraccionaria: con radicales de índices distintos no hay forma
  directa de multiplicar, pero \azul{como potencias es una simple suma de
  fracciones}.}
\end{frame}

\begin{frame}{Cuándo una raíz no existe}
  \scriptsize
  \begin{center}
  \begin{tabular}{@{}p{3.0cm}p{3.2cm}p{5.6cm}@{}}
    \toprule
    \textbf{Caso} & \textbf{Ejemplo} & \textbf{Resultado} \\
    \midrule
    Índice par, radicando positivo & $\sqrt{16}$ &
      $4$ \ \azul{(la raíz principal, positiva)} \\
    \addlinespace
    Índice par, radicando negativo & $\sqrt{-9}$ &
      \destacar{No existe} en los reales \\
    \addlinespace
    Índice impar, radicando positivo & $\sqrt[3]{27}$ & $3$ \\
    \addlinespace
    Índice impar, radicando negativo & $\sqrt[3]{-8}$ &
      $-2$ \ \azul{(sí existe)} \\
    \bottomrule
  \end{tabular}
  \end{center}

  \vspace{-0.02cm}
  \begin{alertblock}{Por qué $\sqrt{-9}$ no existe en $\mathbb{R}$}
    \tiny
    Porque ningún real elevado al cuadrado da negativo: $3^{2}=9$ y
    $(-3)^{2}=9$. \destacar{Es el patrón de la sesión 1}: cuando una operación
    se sale del conjunto, hace falta uno mayor —aquí, los \azul{números
    complejos}, que no son parte de este curso—.
  \end{alertblock}

  \vspace{0.02cm}
  {\tiny\color{uniGris} Ojo con la notación: $\sqrt{16}=4$, \destacar{no}
  $\pm 4$; el símbolo denota la raíz \azul{principal}. El $\pm$ aparece al
  resolver $x^{2}=16$, que es otra cosa (sesión 15).}
\end{frame}

\begin{frame}{Ejercicio 3 — taller principal}
  \begin{alertblock}{En grupos de tres, 20 minutos. Se socializa en el tablero.}
    \scriptsize
    \begin{enumerate}\setlength{\itemsep}{0.0cm}
      \item Simplifique: \ $\sqrt{300}$ \ y \ $\sqrt[3]{250}$
      \item Calcule sin calculadora: \ $8^{2/3}$, \ $16^{3/4}$, \
            $27^{-2/3}$
      \item Opere: \ $\sqrt{50}+\sqrt{18}-\sqrt{8}$
      \item Escriba como un solo radical: \ $\sqrt[3]{a}\cdot\sqrt{a}$
      \item ¿Es cierto que $\sqrt{25+144}=\sqrt{25}+\sqrt{144}$? Compruebe.
      \item El área de un cuadrado es $180\ \text{cm}^{2}$. ¿Cuánto mide el
            lado, en forma exacta y aproximada?
    \end{enumerate}
  \end{alertblock}
\end{frame}

\begin{frame}{Ejercicio 3 — solución}
  \scriptsize
  \begin{columns}[T,onlytextwidth]
    \begin{column}{0.49\textwidth}
      \begin{block}{Puntos 1 a 3}
        \tiny
        \textbf{1.} $\sqrt{2^{2}3\cdot 5^{2}} = \destacar{10\sqrt{3}}$; \
        $\sqrt[3]{2\cdot 5^{3}} = \destacar{5\sqrt[3]{2}}$

        \vspace{0.08cm}
        \textbf{2.} $8^{2/3}=(\sqrt[3]{8})^{2}=\destacar{4}$; \
        $16^{3/4}=2^{3}=\destacar{8}$; \
        $27^{-2/3}=\tfrac{1}{3^{2}}=\destacar{\tfrac{1}{9}}$

        \vspace{0.08cm}
        \textbf{3.} $5\sqrt2+3\sqrt2-2\sqrt2 = \destacar{6\sqrt{2}}$
      \end{block}
    \end{column}
    \begin{column}{0.49\textwidth}
      \begin{block}{Puntos 4 a 6}
        \tiny
        \textbf{4.} $a^{1/3}a^{1/2}=a^{5/6}=\destacar{\sqrt[6]{a^{5}}}$

        \vspace{0.08cm}
        \textbf{5.} \destacar{No}: $\sqrt{169}=13$, pero $5+12=17$.

        \vspace{0.08cm}
        \textbf{6.} $\ell=\sqrt{180}=\destacar{6\sqrt{5}}
        \approx 13{,}42$ cm
      \end{block}
    \end{column}
  \end{columns}

  \vspace{0.08cm}
  {\tiny\color{uniGris} El punto \destacar{6} muestra para qué sirve
  simplificar: $6\sqrt{5}$ es la respuesta \azul{exacta} y $13{,}42$ es una
  aproximación. En un examen, dejar $\sqrt{180}$ sin simplificar suele costar
  medio punto.}
\end{frame}

\begin{frame}{Errores frecuentes en esta sesión}
  \scriptsize
  \begin{center}
  \begin{tabular}{@{}p{5.0cm}p{7.0cm}@{}}
    \toprule
    \textbf{Error} & \textbf{Lo correcto} \\
    \midrule
    $\sqrt{a+b}=\sqrt{a}+\sqrt{b}$ &
      No existe esa regla; solo vale para $\cdot$ y $\div$ \\
    \addlinespace
    $\sqrt{12}+\sqrt{27}=\sqrt{39}$ &
      Simplifique primero: $2\sqrt3+3\sqrt3=5\sqrt3$ \\
    \addlinespace
    $\sqrt{16}=\pm 4$ &
      La raíz es la principal: $4$ \\
    \addlinespace
    $\sqrt[3]{-8}$ ``no existe'' &
      Con índice impar sí existe: $-2$ \\
    \addlinespace
    $8^{2/3}=\sqrt{8^{3}}$ &
      El $3$ es el índice: $\sqrt[3]{8^{2}}=4$ \\
    \addlinespace
    Dejar $\sqrt{180}$ como respuesta &
      Extraiga factores: $6\sqrt{5}$ \\
    \bottomrule
  \end{tabular}
  \end{center}
\end{frame}

\begin{frame}{Cierre de la sesión}
  \scriptsize
  \begin{block}{Lo que hay que llevarse}
    \tiny
    \textbf{1.} $\sqrt[n]{a^{m}} = a^{m/n}$: índice abajo, exponente
    arriba. \quad
    \textbf{2.} La raíz se reparte en productos y cocientes,
    \destacar{nunca en sumas}. \quad
    \textbf{3.} Para simplificar, agrupe en bloques del tamaño del índice.
    \quad
    \textbf{4.} Solo se suman radicales semejantes, y hay que simplificar
    antes de decidirlo. \quad
    \textbf{5.} Índice par con radicando negativo: no existe en
    $\mathbb{R}$.
  \end{block}

  \vspace{0.1cm}
  \begin{alertblock}{Para la próxima sesión}
    \tiny
    Resolver los \destacar{20 ejercicios propuestos} y la guía de la unidad en
    Aulas Virtuales.

    \vspace{0.05cm}
    En la sesión 9 vemos \azul{racionalización}: cómo quitar los radicales del
    denominador. Se apoya por completo en lo de hoy y en los productos
    notables, que veremos formalmente en la Unidad 3 pero que allí usaremos
    ya.
  \end{alertblock}
\end{frame}

\begin{frame}{Ejercicios propuestos}
  \vspace{-0.1cm}
  \begin{beamercolorbox}[wd=\textwidth,sep=0.10cm]{block body}
    \centering\scriptsize
    \textbf{1--7:} simplifique. \quad \textbf{8--12:} calcule sin
    calculadora. \quad \textbf{13--17:} opere. \quad
    \textbf{18--20:} aplicación.
  \end{beamercolorbox}
  \listatiny
  \vspace{0.15cm}
  \begin{columns}[T,onlytextwidth]
    \begin{column}{0.47\textwidth}
      \begin{enumerate}\setlength{\itemsep}{0.16cm}
        \item $\sqrt{98}$
        \item $\sqrt{242}$
        \item $\sqrt[3]{135}$
        \item $\sqrt[4]{162}$
        \item $\sqrt{45x^{2}}$
        \item $\sqrt{32x^{5}y^{4}}$
        \item $\sqrt[3]{81a^{6}}$
        \item $25^{3/2}$
        \item $64^{2/3}$
        \item $81^{-3/4}$
      \end{enumerate}
    \end{column}
    \begin{column}{0.50\textwidth}
      \begin{enumerate}\setlength{\itemsep}{0.16cm}
        \setcounter{enumi}{10}
        \item $\left(\sqrt[5]{3}\right)^{10}$
        \item $\sqrt[3]{-125}$
        \item $\sqrt{7}\cdot\sqrt{28}$
        \item $\dfrac{\sqrt{150}}{\sqrt{6}}$
        \item $\sqrt{75}-\sqrt{12}$
        \item $4\sqrt{18}+2\sqrt{50}$
        \item $\sqrt[4]{a}\cdot\sqrt{a}$ como un solo radical
        \item ¿Existe $\sqrt[4]{-16}$? Justifique
        \item El área de un cuadrado es $128$ m$^2$: halle el lado exacto
        \item Un cubo tiene volumen $54$ cm$^3$: halle la arista exacta
      \end{enumerate}
    \end{column}
  \end{columns}
\end{frame}

\begin{frame}{Respuestas}{Autoevaluación de los ejercicios propuestos}
  \listatiny
  \vspace{0.15cm}
  \begin{columns}[T,onlytextwidth]
    \begin{column}{0.47\textwidth}
      \begin{enumerate}\setlength{\itemsep}{0.16cm}
        \item $7\sqrt{2}$
        \item $11\sqrt{2}$
        \item $3\sqrt[3]{5}$
        \item $3\sqrt[4]{2}$
        \item $3x\sqrt{5}$
        \item $4x^{2}y^{2}\sqrt{2x}$
        \item $3a^{2}\sqrt[3]{3}$
        \item $125$
        \item $16$
        \item $\tfrac{1}{27}$
      \end{enumerate}
    \end{column}
    \begin{column}{0.50\textwidth}
      \begin{enumerate}\setlength{\itemsep}{0.16cm}
        \setcounter{enumi}{10}
        \item $9$
        \item $-5$
        \item $14$
        \item $5$
        \item $3\sqrt{3}$
        \item $22\sqrt{2}$
        \item $a^{3/4}=\sqrt[4]{a^{3}}$
        \item No: índice par y radicando negativo
        \item $8\sqrt{2}\approx 11{,}31$ m
        \item $3\sqrt[3]{2}\approx 3{,}78$ cm
      \end{enumerate}
    \end{column}
  \end{columns}
\end{frame}

\begin{frame}{Referencias}
  \footnotesize
  \begin{itemize}
    \item[{[1]}] Ramos del Olmo, Á. M. (2016). \textit{Matemáticas básicas
      para el acceso a la universidad}. Ediciones Pirámide.
      \texttt{https://elibro.net/es/ereader/uniminuto/49134}
    \item[{[2]}] González Pulido, J. W., Chacón Benavides, J. A. \&
      Fonseca Correa, L. Á. (2023). \textit{Matemática básica}
      (1.\textsuperscript{a} ed.). Universidad Mariana.
      \texttt{https://elibro.net/es/ereader/uniminuto/280647}
    \item[{[3]}] Torrijos, M. (2018). \textit{Apuntes de matemática básica}.
      Universidad Católica de Colombia.
      \texttt{https://elibro.net/es/ereader/uniminuto/197048}
    \item[{[4]}] Ramírez, A. \& Rojas, L. (2016). \textit{Matemáticas básicas
      con aplicaciones a la ingeniería}. Ecoe Ediciones.
      \texttt{https://elibro.net/es/ereader/uniminuto/114355}
  \end{itemize}

  \vspace{0.3cm}
  \begin{center}
    {\Large\color{uniAzul}\textbf{¿Preguntas?}}
  \end{center}
\end{frame}

\end{document}
